\chapter{مقدمه}

این فصل تصویر کلی پروژه را توضیح می‌دهد: چرا چنین سامانه‌ای لازم است، کاربر با آن چه می‌کند، و گزارش از چه زاویه‌ای به پیاده‌سازی نگاه می‌کند. متن تلاش می‌کند فنی باشد، اما خشک و بریده‌بریده نباشد؛ یعنی خواننده ابتدا حس محصول را بگیرد و بعد وارد جزئیات معماری شود.

\section{زمینه پروژه}
\lr{IranAPI Hub} یک پلتفرم فارسی برای معرفی، جستجو، انتشار و مصرف \lr{API} است. فضای پروژه به بازارچه‌های \lr{API} نزدیک است؛ جایی که کاربر می‌تواند سرویس‌ها را ببیند، مستندات را بخواند، پلن‌ها را مقایسه کند و بعد از ورود، مصرف و دسترسی‌های خود را مدیریت کند. این نگاه با مدل رایج \lr{API Hub} و بازارچه‌های \lr{API} هم‌خوان است \cite{rapidapi}.

از نظر فنی، پروژه روی ترکیب \lr{Python}، \lr{Django REST Framework}، \lr{React}، \lr{Vite}، \lr{MongoDB} و \lr{Docker Compose} شکل گرفته است \cite{python,django,drf,react,vite,mongodb,docker}. این ترکیب برای سامانه‌ای مناسب است که هم به یک \lr{API} ساخت‌یافته نیاز دارد و هم به یک رابط کاربری سریع و قابل توسعه.

\section{بیان مسئله}
کاربر چنین سامانه‌ای معمولاً نمی‌خواهد فقط یک فهرست ساده از سرویس‌ها ببیند. او باید بتواند \lr{API}ها را جستجو کند، دسته‌بندی‌ها را بررسی کند، مستندات هر \lr{Endpoint} را بخواند، قیمت یا پلن را بسنجد و بعد از ورود، وضعیت مصرف، کلیدها، اشتراک و سازمان خود را مدیریت کند. توسعه‌دهنده هم نیاز دارد \lr{API} جدید منتشر کند، نمونه درخواست بفرستد و جریان کاری خود را در \lr{Studio} تنظیم کند.

مسئله اصلی، پراکندگی تجربه کاربر است. وقتی مستندات، پرداخت، احراز هویت، مصرف و ابزار تست در چند مسیر جدا باشند، کاربر برای رسیدن به نتیجه باید چند بار زمینه ذهنی خود را عوض کند. این پروژه تلاش می‌کند همین اجزا را در یک تجربه منسجم جمع کند؛ تجربه‌ای که هم برای کاربر عمومی قابل فهم باشد و هم برای توسعه‌دهنده کاربرد عملی داشته باشد.

\section{اهداف}
اهداف اصلی پروژه عبارت‌اند از:
\begin{itemize}
  \item ارائه کاتالوگ \lr{API} با جستجو، دسته‌بندی، امتیاز، پلن و مستندات.
  \item فراهم‌کردن ثبت‌نام، ورود، خروج و نشست کاربر با مدل رایج احراز هویت وب.
  \item نمایش \lr{Dashboard} توسعه‌دهنده شامل پروفایل، \lr{API Key}، مصرف، سازمان و اشتراک.
  \item پشتیبانی از انتشار \lr{API}، اجرای نمونه درخواست در \lr{Caller} و ساخت جریان در \lr{Studio}.
  \item ارائه قرارداد \lr{OpenAPI} برای خوانایی بهتر مسیرها و امکان اتصال ابزارهای دیگر \cite{openapi}.
\end{itemize}

این اهداف فقط به ظاهر سامانه مربوط نیستند. کاتالوگ و مستندات بخش عمومی محصول را می‌سازند؛ ورود، اشتراک و مصرف بخش خصوصی را شکل می‌دهند؛ و \lr{OpenAPI} کمک می‌کند قرارداد سرویس برای انسان و ابزار قابل خواندن باشد.

\section{دامنه و محدودیت‌ها}
دامنه گزارش، نسخه فعلی پروژه و قابلیت‌هایی است که در آن قابل بررسی‌اند: وب‌اپلیکیشن فارسی، \lr{API} بک‌اند، اجرای محلی و جریان‌های اصلی کاربر. اجرای محلی با \lr{Docker Compose} در همین راستا مطرح می‌شود، چون برای بالا آوردن چند سرویس کنار هم طراحی شده است \cite{docker}.

در کنار قابلیت‌ها، چند محدودیت هم باید روشن گفته شود. تنظیمات تولیدی هنوز جای کار دارد، بخشی از رفتارها برای محیط توسعه مناسب‌ترند، پرداخت واقعی و اتصال کامل به سرویس‌های بیرونی نیاز به سخت‌سازی دارد، و آزمون مرورگر باید در چرخه ارزیابی منظم‌تر شود. بنابراین گزارش، پروژه را به عنوان یک نمونه کامل و قابل توسعه بررسی می‌کند، نه یک محصول تولیدی تمام‌شده.

\section{ذی‌نفعان}
ذی‌نفعان اصلی پروژه شامل کاربر مهمان، توسعه‌دهنده احرازشده، مدیر سیستم، ارائه‌دهنده ورود اجتماعی، سرویس مقصد \lr{API} و بازارهای بیرونی مانند \lr{RapidAPI} هستند \cite{rapidapi}. هر کدام از این گروه‌ها انتظار متفاوتی دارند: مهمان به شناخت سریع سرویس نیاز دارد، توسعه‌دهنده به ابزار و داده خصوصی، مدیر به کنترل سامانه، و سرویس‌های بیرونی به قرارداد قابل اعتماد.

\section{سازمان‌دهی گزارش}
فصل دوم فناوری‌ها را معرفی می‌کند. فصل سوم نیازمندی‌ها را دسته‌بندی می‌کند. فصل چهارم معماری را توضیح می‌دهد. فصل پنجم به پیاده‌سازی می‌پردازد. فصل ششم آزمون‌ها را گزارش می‌کند. فصل هفتم نتایج و چالش‌ها را جمع می‌کند و فصل هشتم مسیر ادامه کار را نشان می‌دهد.

\section{جمع‌بندی فصل}
در این فصل مشخص شد که پروژه یک پلتفرم اشتراک‌گذاری و مصرف \lr{API} است؛ نه فقط یک تمرین بک‌اند یا یک رابط کاربری ساده. همین ماهیت دوگانه باعث می‌شود گزارش هم به تجربه کاربر توجه کند و هم به قرارداد فنی بین فرانت‌اند، بک‌اند و داده.
